\section{Background and Related Work}
\label{sec:related}

\paragraph{The \texttt{dot} layout pipeline.}
GraphViz's \texttt{dot} engine draws directed graphs using the layered approach
introduced by Sugiyama et al.~\cite{sugiyama1981methods} and refined by Gansner
et al.~\cite{gansner1993technique}. Layout proceeds in four phases, each
depending on the last: \emph{ranking} assigns nodes to layers, \emph{crossing
minimization} orders nodes within each layer to reduce edge crossings,
\emph{coordinate assignment} fixes $x$ positions, and \emph{spline routing}
draws the edges. Ranking and coordinate assignment are both solved with network
simplex; crossing minimization is an iterative heuristic. GraphViz exposes a
\texttt{-Gphase=N} attribute that runs a prefix of this pipeline and stops,
which we use for phase attribution (Section~\ref{sec:method}).

\paragraph{Agents modifying real code.}
Benchmarks such as SWE-bench~\cite{jimenez2024swebench} evaluate agents on
resolving real repository issues, where a test suite provides the success
signal. Performance work differs in that no test can tell you whether an agent
found the \emph{right} thing to optimize, and a change that is correct but
useless passes every test. Prior work on learned compiler optimization
\cite{cummins2023llmcompiler} targets transformations chosen in advance; here
the choice of target is the object of study.

\paragraph{Prior work on this codebase.}
An earlier study by some of the present authors applied a multi-model ensemble
to the same GraphViz layout code and reported a $3.78\times$ speedup from
OpenMP parallelization of crossing minimization~\cite{estrada2025openmp}. That
work has not been peer-reviewed. Our measurements do not reproduce it, and
Section~\ref{sec:discussion} gives the phase composition and Amdahl ceiling
that bound what parallelizing that phase can yield on this workload. We note
also that its reported prompt template specifies the target speedup as an input
requirement, which makes the reported figure difficult to interpret as an
independent measurement. The present paper shares only the codebase and the
general question; its protocol, measurements, and conclusions are separate.
